The two models fail in opposite ways. MedGemma, trained on medical images, is clearly better at reading the lesion: it is more accurate, more stable when the background changes, and its answers about dermoscopic structures agree better with dermatologists. But it does not recognise when there is nothing to read. With the lesion gone it still gives a confident answer for half the images, and across all images it is roughly twice as confident as it is right. Qwen2.5-VL is the opposite: it is cautious, often to the point of declining to answer, and it almost always notices that the lesion is missing, but its diagnoses change for a quarter of images when only the surrounding skin is edited, and its feature answers are close to ``yes to everything''.

None of this is visible in an accuracy number. A user who saw only that MedGemma reaches 43\% balanced accuracy on smartphone photos would not learn that, on photos with the lesion painted out, its average confidence in its top diagnosis is still 70\%. For the use that motivates this work, a referral aid for health workers far from a dermatologist, that is the failure that matters most: a confident wrong answer is worse than ``I cannot tell''. The lesion-removal and background tests cost a few lines of image processing and one extra model call per image, and they could be reported alongside accuracy for any medical VLM.

Some results are smaller than expected. Earlier work found that surgical-marker ink pushed a convolutional classifier towards melanoma \cite{winkler2019}; here marker lines changed diagnoses but did not move the melanoma probability of either model measurably, while a ruler did for MedGemma. The skin-tone edit is crude, so the effect of darker skin on Qwen2.5-VL should be checked with real images of darker skin before drawing any conclusion about fairness.

The pilot suggests three directions for a full study: (1) training or prompting methods that let a model abstain when the lesion is missing or unclear, evaluated with exactly these counterfactual tests; (2) grounding the model's stated features in the image, for example by asking it to point to each structure and checking the location against a segmentation; and (3) repeating the tests on larger and fine-tuned models and on images from South Asian patients, where none of the public datasets used here applies.
